\section{Metabolic Adaptations in Developing Neocortical Cell Types}

\subsection{Are Progenitors OXPHOS-ready?}
We further characterized metabolic differences across cell types in the developing neocortex at mid-neurogenesis (E14.5), focusing on how oxygen availability influences their bioenergetic strategies. Our findings suggest that oxygen acts as a limiting factor for OXPHOS, particularly in progenitor cells. Gene expression data together with both, E-Flux and weighted iMAT indicate that progenitors possess the molecular machinery for OXPHOS, as evidenced by increasing CV activity with rising oxygen levels. Moreover, TCA cycle activity persists and is even increased under limited oxygen conditions in NPCs, as indicated by sustained flux through SUCOAS.
\\

These observations raise the questions of whether progenitor cells are intrinsically equipped for OXPHOS and are primarily limited by oxygen availability, rather than by the absence of metabolic machinery. 
\\

From a developmental perspective, this hypothesis is consistent with the spatial and temporal dynamics of neocortical vascularization. Angiogenesis begins in parallel with neurogenesis, leading to the formation of distinct vascular niches within the developing cortex. A relatively hypoxic environment is maintained near the VZ, whereas more oxygen-rich, stable perivascular niches emerge in the SVZ \citep{komabayashi-suzukiSpatiotemporallyDependentVascularization2019}.
Given the different associations of APs and BPs within the vascular niches, it would not be surprising that they possess different metabolic profiles. 
\\

Experimental evidence further supports the idea that progenitors possess an already functional mitochondrial machinery. For example, inhibition of  CI impaired proliferation and terminal differentiation of NPCs \citep{CABELLO-RIVERAMitochondrialComplexFunction2019,vandenameeleNeuralStemCell2019}.
\\

Furthermore, progenitors are clearly capable of engaging OXPHOS when conditions permit. 
\textit{In vitro} studies demonstrate that NPCs cultured under 3\% vs. 21\% oxygen  exhibit marked differences in mitochondrial properties. Under low oxygen, NPCs display reduced mitochondrial mass, yet maintain energy production through an increased mitochondrial membrane potential compared to cells grown in normoxia \citep{lagesLowOxygenAlters2015}.
This suggests a more efficient use of the electron transport chain under oxygen-limited conditions rather than its absence.
\\

Taken together, these findings support a model in which progenitor cells are metabolically flexible. While NPCs possess a functional electron transport chain, they predominantly rely on glycolysis. Whether this preference is primarily driven by environmental constraints, such as oxygen availability, or reflects an intrinsic requirement of glycolysis to maintain progenitor identity and stemness, remains an open question. 


\section{Coupling of Oxygen Availability and Cell Cycle-Dependent Metabolic States in AP}
Such metabolic flexibility is likely advantageous during cortical development, where progenitors experience spatial and temporal variations in oxygen availability. In this context, the ability to modulate OXPHOS in response to oxygen levels may not only reflect environmental constraints but also support dynamic cellular processes. 
\\

One such process is INM, in which AP nuclei move along the apical-basal axis in synchrony with the cell cycle \citep{grosseCellDynamicsFetal2011}.
Given that different regions of the developing cortex are associated with distinct oxygen levels, this movement could expose cells to fluctuating metabolic conditions over the course of a single cell cycle. Concomitantly, metabolism itself is known to vary across cell cycle phases, with shifts between glycolytic and oxidative states reported in proliferating cells \citep{zylstraMetabolicDynamicsCell2022,liuSkp2DictatesCell2021}.
\\


In this light, the oxygen-dependent metabolic flexibility observed in our data may enable APs to adapt their energy metabolism both to local oxygen availability and to phase-specific metabolic demands during the cell cycle. While this remains speculative, it provides a potential explanation for why progenitors retain a functional electron transport chain despite predominantly relying on glycolysis. 

\section{Cell-Type specific Metabolic Specialization Revealed by Context-Specific GEMs}
\subsection{Basal Progenitors exhibit distinct Fatty Acid Oxidation and Serine Utilization Programs}
BP  exhibited several unique reactions associated with fatty acid metabolism, including multiple acyl-CoA dehydrogenase reactions, which catalyze the first step of $\beta$-oxidation by reducing fatty acids while transferring electrons to ubiquinone, generating ubiquinol that subsequently feeds electrons into the ETC. Together, these reactions suggest active FAO in BPs when oxygen is available, consistent with previous studies reporting substantial $\beta$-oxidation activity in NPCs from the SVZ \citep{STOLLNeuralStemCells2015}.
\\

This raises the possibility that BPs may not always have consistent access to glucose compared to neurons in deep or  upper cortical layers. If glucose is scarce, BPs could prioritize FAO as an alternative energy source. In fact, it might be more efficient to oxidize glucose to lactate, which does not consume oxygen, and reserve the limited oxygen for $\beta$-oxidation of fatty acids, maximizing ATP yield under oxygen-restricted conditions.
\\

Beside reactions involved in FAO, reactions concerning serine metabolism were found. Serine, as a glucogenic amino acid, can be converted into pyruvate via the  enzyme serine-ammona-lyase, which we identified as a unique and irreversible reaction for BPs (r0060). This pathway additionally provides a route to bypass glucose dependency, allowing BPs to use serine as an alternative carbon source for lactate production or entry into the TCA cycle. Furthermore, a serine-rich diet has been shown to enhance proliferation of NSCs 
\textit{in vitro} and \textit{in vivo}, increasing cell proliferation and density of NSCs as well as NPCs \citep{sultanDserineIncreasesAdult2013,than-trongLserineDietRestores2026,kawakamiImpairedNeurogenesisEmbryonic2009}.
Defects in serine production due to PHGDH loss-of-function lead to embryonic lethality with severe morphological defects in the CNS, including incomplete cerebellar development \citep{yoshidaTargetedDisruptionMouse2004}.
While the exact mechanisms underlying these phenotypes remain to be uncovered, they underscore the role of serine metabolism during neurogenesis. In this context, the conversion of serine into pyruvate may represent a mechanism to locally provide pyruvate for pathways such as TCA cycle, biosynthetic metabolism or lactate production.


\subsection{Shared Glutamine-Glutamate Metabolism links Apical Progenitors and Deep-Layer Neurons}
Moreover, we identified reactions shared across all cycling AP populations and both DLN populations that were associated with glutamine and alanine exchange (r1434, ALAt2r, EX\_ala\_L\_e, GLNtm), as well as glutamate production from glutamine via mitochondrial glutaminase (GLUNm), an irreversible reaction. In addition, reversible alanine aminotransferase reactions in both the cytosol and mitochondria (ALATA\_L and r0081, respectively) interconvert alanine and $\alpha$KG to form glutamate and pyruvate.  Together, these reactions connect glutaminolysis, a metabolic pathway in which glutamine is first converted into glutamate and subsequently deaminated to  $\alpha$KG, to the TCA cycle.
\\

In comparison to BPs, active glutaminolysis has specifically been described in mural  APs, where pharmacological inhibition reduces progenitor proliferation during neurogenesis \citep{nambaHumanSpecificARHGAP11BActs2020,journiacCellMetabolicAlterations2020}.
Increased glutaminolysis has further been reported in both mouse and human APs and has been proposed as a hallmark of proliferative embryonic stem cell states \citep{yangGlutaminolysisHallmarkCancer2017,journiacCellMetabolicAlterations2020}. Through the production of aKG, glutaminolysis can support TCA cycle activity independently of glycolysis.
\\

\textit{Namba et al.} proposed that APs may operate a ‘¾ TCA cycle’, in which the cycle terminates at OAA rather than continuing through citrate synthesis, thereby redirecting intermediates toward anabolic pathways such as amino acid, nucleotide, and lipid synthesis \citep{nambaMetabolicRegulationNeocortical2021}.
In addition, $\alpha$KG itself has been implicated in epigenetic regulation and stem cell differentiation \citep{yaoEpigeneticMechanismsNeurogenesis2016}.
\\

Since glutaminolysis requires a continuous glutamine supply, the identified glutamine transport reaction (GLNtm) may indicate glutamine uptake. Potential glutamine sources for APs include the cerebrospinal fluid, given their proximity to the lateral ventricles, as well as the developing vasculature. For DLNs, uptake from the bloodstream may additionally be possible during early fetal stages before full maturation of the blood-brain barrier \citep{goasdoueReviewBloodbrainBarrier2017}.
\\


Notably, these reactions were shared specifically between APs and DLNs, potentially reflecting metabolic programs associated with their lineage relationship. However, the shared pathways may also fulfill distinct functions in the two cell types. In APs, glutaminolysis appears to primarily support proliferative and biosynthetic demands during neurogenesis. In contrast, DLNs may increasingly utilize glutamate for neuron-specific functions. Besides serving as a metabolic substrate, glutamate is the principal excitatory neurotransmitter of the mature CNS, with glutaminase representing the major glutamate-producing enzyme in the brain 
\citep{marquezNewInsightsBrain2009}. Furthermore, glutamate and glutamine can serve as alternative energy substrates for neurons under hypoglycemic conditions \citep{mckennaGlutamateglutamineCycleNot2007}. Together, these findings suggest that glutamine-glutamate metabolism may support both proliferative progenitor programs and early neuronal metabolic adaptation during corticogenesis. 

\section{Conclusion and Future Perspectives}
Future studies should take advantage of these context-specific GEMs to simulate reaction fluxes and further refine the models by accounting for nutrient availability. While our collaborators provided untargeted metabolomic data, many metabolites remained unidentified or were not included in our model. Incorporating targeted metabolomic data or literature-derived metabolite concentrations could improve model accuracy. Additionally, comparing cell types across multiple developmental time points, both earlier and later than E14.5, would reveal temporal metabolic shifts and may help define clearer metabolic profiles. Finally, experimental validation of model predictions would strengthen the biological relevance of these findings. 
